\documentclass{amsart}
% For dealing with references we use the comment environment
\usepackage{verbatim}
\newenvironment{reference}{\comment}{\endcomment}
%\newenvironment{reference}{}{}
\newenvironment{slogan}{\comment}{\endcomment}
\newenvironment{history}{\comment}{\endcomment}

% For commutative diagrams we use Xy-pic
\usepackage[all]{xy}

% We use 2cell for 2-commutative diagrams.
\xyoption{2cell}
\UseAllTwocells

% We use multicol for the list of chapters between chapters
\usepackage{multicol}

% This is generally recommended for better output
\usepackage{lmodern}
\usepackage[T1]{fontenc}

% For cross-file-references

% Package for hypertext links:
\usepackage{hyperref}

% For any local file, say "hello.tex" you want to link to please

% Theorem environments.
%
\theoremstyle{plain}
\newtheorem{theorem}[subsection]{Theorem}
\newtheorem{proposition}[subsection]{Proposition}
\newtheorem{lemma}[subsection]{Lemma}

\theoremstyle{definition}
\newtheorem{definition}[subsection]{Definition}
\newtheorem{example}[subsection]{Example}
\newtheorem{exercise}[subsection]{Exercise}
\newtheorem{situation}[subsection]{Situation}

\theoremstyle{remark}
\newtheorem{remark}[subsection]{Remark}
\newtheorem{remarks}[subsection]{Remarks}

\numberwithin{equation}{subsection}

% Macros
%
\def\lim{\mathop{\mathrm{lim}}\nolimits}
\def\colim{\mathop{\mathrm{colim}}\nolimits}
\def\Spec{\mathop{\mathrm{Spec}}}
\def\Hom{\mathop{\mathrm{Hom}}\nolimits}
\def\Ext{\mathop{\mathrm{Ext}}\nolimits}
\def\SheafHom{\mathop{\mathcal{H}\!\mathit{om}}\nolimits}
\def\SheafExt{\mathop{\mathcal{E}\!\mathit{xt}}\nolimits}
\def\Sch{\mathit{Sch}}
\def\Mor{\mathop{\mathrm{Mor}}\nolimits}
\def\Ob{\mathop{\mathrm{Ob}}\nolimits}
\def\Sh{\mathop{\mathit{Sh}}\nolimits}
\def\NL{\mathop{N\!L}\nolimits}
\def\CH{\mathop{\mathrm{CH}}\nolimits}
\def\proetale{{pro\text{-}\acute{e}tale}}
\def\etale{{\acute{e}tale}}
\def\QCoh{\mathit{QCoh}}
\def\Ker{\mathop{\mathrm{Ker}}}
\def\Im{\mathop{\mathrm{Im}}}
\def\Coker{\mathop{\mathrm{Coker}}}
\def\Coim{\mathop{\mathrm{Coim}}}

% Boxtimes
%
\DeclareMathSymbol{\boxtimes}{\mathbin}{AMSa}{"02}

%
% Macros for moduli stacks/spaces
%
\def\QCohstack{\mathcal{QC}\!\mathit{oh}}
\def\Cohstack{\mathcal{C}\!\mathit{oh}}
\def\Spacesstack{\mathcal{S}\!\mathit{paces}}
\def\Quotfunctor{\mathrm{Quot}}
\def\Hilbfunctor{\mathrm{Hilb}}
\def\Curvesstack{\mathcal{C}\!\mathit{urves}}
\def\Polarizedstack{\mathcal{P}\!\mathit{olarized}}
\def\Complexesstack{\mathcal{C}\!\mathit{omplexes}}
% \Pic is the operator that assigns to X its picard group, usage \Pic(X)
% \Picardstack_{X/B} denotes the Picard stack of X over B
% \Picardfunctor_{X/B} denotes the Picard functor of X over B
\def\Pic{\mathop{\mathrm{Pic}}\nolimits}
\def\Picardstack{\mathcal{P}\!\mathit{ic}}
\def\Picardfunctor{\mathrm{Pic}}
\def\Deformationcategory{\mathcal{D}\!\mathit{ef}}

\setlength{\emergencystretch}{3em}
\begin{document}
\title[Stacks Project, Tag 0BRM]{421 --- Abhyankar lemma for tame extensions of discrete valuation rings}
\maketitle
\noindent Copyright (C) 2005--2025 Johan de Jong. Stacks Project authors. Source: \href{https://stacks.math.columbia.edu/tag/0BRM}{Tag 0BRM}.

\noindent Editorial difficulty note (not author text). Difficulty H3: The proof first kills a unit factor through a finite etale extension and descends the desired formal smoothness. It then reduces to adjoining a root of a uniformizer, computes ramification indices and residue separability after normalization, and proves that the new local extensions are unramified..

\noindent Editorial provenance note (not author text). History: excerpt from revision \texttt{a0444\allowbreak 6e57e\allowbreak c1fbc\allowbreak 252a8\allowbreak 71afc\allowbreak ec775\allowbreak 2fb28\allowbreak 07b14}, more-algebra.tex, lines 34036--34135. Reference presentation was adapted; original-author context and core support blocks were retained or added. The mathematical wording is unchanged. Licensed under GNU FDL 1.2 or later; no invariant sections or cover texts. See ../licenses/COPYING. Context blocks reproduce author definitions and supporting results; they are not additional counted problems.

\begin{remark}
\label{remark-construction}
Let $A \to B$ be an extension of discrete valuation rings with fraction
fields $K \subset L$. Let $K_1/K$ be a finite extension of
fields. Let $A_1 \subset K_1$ be the integral closure of $A$ in $K_1$.
On the other hand, let $L_1 = (L \otimes_K K_1)_{red}$. Then $L_1$ is a
nonempty finite product of finite field extensions of $L$. Let $B_1$ be
the integral closure of $B$ in $L_1$. We obtain compatible commutative
diagrams
$$
\vcenter{
\xymatrix{
L \ar[r] & L_1 \\
K \ar[u] \ar[r] & K_1 \ar[u]
}
}
\quad\text{and}\quad
\vcenter{
\xymatrix{
B \ar[r] & B_1 \\
A \ar[u] \ar[r] & A_1 \ar[u]
}
}
$$
In this situation we have the following
\begin{enumerate}
\item By Algebra, Lemma \href{https://stacks.math.columbia.edu/tag/09IG}{lemma integral closure Dedekind (Tag 09IG)}
the ring $A_1$ is a Dedekind domain and $B_1$ is a finite product of
Dedekind domains.
\item Note that $L \otimes_K K_1 = (B \otimes_A A_1)_\pi$ where $\pi \in A$
is a uniformizer and that $\pi$ is a nonzerodivisor on $B \otimes_A A_1$. 
Thus the ring map $B \otimes_A A_1 \to B_1$ is integral with kernel
consisting of nilpotent elements. Hence $\Spec(B_1) \to \Spec(B \otimes_A A_1)$
is surjective on spectra
(Algebra, Lemma \href{https://stacks.math.columbia.edu/tag/00GQ}{lemma integral overring surjective (Tag 00GQ)}).
The map $\Spec(B \otimes_A A_1) \to \Spec(A_1)$ is surjective as
$A_1/\mathfrak m_A A_1 \to
B/\mathfrak m_AB \otimes_{\kappa_A} A_1/\mathfrak m_A A_1$
is an injective ring map with $A_1/\mathfrak m_A A_1$ Artinian.
We conclude that $\Spec(B_1) \to \Spec(A_1)$ is surjective.
\item Let $\mathfrak m_i$, $i = 1, \ldots n$ with $n \geq 1$ be the
maximal ideals of $A_1$. For each $i = 1, \ldots, n$ let
$\mathfrak m_{ij}$, $j = 1, \ldots, m_i$ with $m_i \geq 1$
be the maximal ideals of $B_1$ lying over $\mathfrak m_i$. We obtain diagrams
$$
\xymatrix{
B \ar[r] & (B_1)_{\mathfrak m_{ij}} \\
A \ar[u] \ar[r] & (A_1)_{\mathfrak m_i} \ar[u]
}
$$
of extensions of discrete valuation rings.
\item If $A$ is henselian (for example complete), then $A_1$ is a
discrete valuation ring, i.e., $n = 1$.
Namely, $A_1$ is a union of finite extensions of $A$ which are domains,
hence local by Algebra, Lemma \href{https://stacks.math.columbia.edu/tag/04GH}{lemma finite over henselian (Tag 04GH)}.
\item If $B$ is henselian (for example complete), then $B_1$
is a product of discrete valuation rings, i.e., $m_i = 1$ for
$i = 1, \ldots, n$.
\item If $K \subset K_1$ is purely inseparable, then $A_1$ and $B_1$
are both discrete valuation rings, i.e., $n = 1$ and $m_1 = 1$.
This is true because for every $b \in B_1$ a $p$-power power of $b$
is in $B$, hence $B_1$ can only have one maximal ideal.
\item If $K \subset K_1$ is finite separable, then $L_1 = L \otimes_K K_1$
and is a finite product of finite separable extensions too. Hence
$A \subset A_1$ and $B \subset B_1$ are finite by
Algebra, Lemma
\href{https://stacks.math.columbia.edu/tag/032L}{lemma Noetherian normal domain finite separable extension (Tag 032L)}.
\item If $A$ is Nagata, then $A \subset A_1$ is finite.
\item If $B$ is Nagata, then $B \subset B_1$ is finite.
\end{enumerate}
\end{remark}


\begin{lemma}[Abhyankar's lemma]
\label{lemma-abhyankar}
Let $A \subset B$ be an extension of discrete valuation rings.
Assume that either the residue characteristic of $A$ is $0$
or it is $p$, the ramification index $e$ is prime to $p$, and
$\kappa_B/\kappa_A$ is a separable field extension.
Let $K_1/K$ be a finite extension. Using the notation of
Remark \ref{remark-construction}
assume $e$ divides the ramification index of $A \subset (A_1)_{\mathfrak m_i}$
for some $i$. Then $(A_1)_{\mathfrak m_i} \subset (B_1)_{\mathfrak m_{ij}}$
is formally smooth in the $\mathfrak m_{ij}$-adic topology
for all $j = 1, \ldots, m_i$.
\end{lemma}

\begin{proof}
Let $\pi \in A$ be a uniformizer. Let $\pi_1$ be a uniformizer
of $(A_1)_{\mathfrak m_i}$. Write $\pi = u \pi_1^{e_1}$ with $u$ a unit
of $(A_1)_{\mathfrak m_i}$ and $e_1$ the ramification index of
$A \subset (A_1)_{\mathfrak m_i}$.

\medskip\noindent
Claim: we may assume that $u$ is an $e$th power in $K_1$.
Namely, let $K_2$ be an extension of $K_1$ obtained by
adjoining a root of $x^e = u$; thus $K_2$ is a factor
of $K_1[x]/(x^e - u)$. Then $K_2/K_1$ is a finite
separable extension (by our assumption on $e$)
and hence $A_1 \subset A_2$ is finite.
Since $(A_1)_{\mathfrak m_i} \to (A_1)_{\mathfrak m_i}[x]/(x^e - u)$
is finite \'etale
(as $e$ is prime to the residue characteristic and $u$ a unit)
we conclude that $(A_2)_{\mathfrak m_i}$ is a factor of
a finite \'etale extension of $(A_1)_{\mathfrak m_i}$ hence
finite \'etale over $(A_1)_{\mathfrak m_i}$ itself.
The same reasoning shows that $B_1 \subset B_2$ induces
finite \'etale extensions
$(B_1)_{\mathfrak m_{ij}} \subset (B_2)_{\mathfrak m_{ij}}$.
Pick a maximal ideal $\mathfrak m'_{ij} \subset B_2$
lying over $\mathfrak m_{ij} \subset B_1$
(of course there may be more than one) and consider
$$
\xymatrix{
(B_1)_{\mathfrak m_{ij}} \ar[r] & (B_2)_{\mathfrak m'_{ij}} \\
(A_1)_{\mathfrak m_i} \ar[u] \ar[r] &
(A_2)_{\mathfrak m'_i} \ar[u]
}
$$
where $\mathfrak m'_i \subset A_2$ is the image.
Now the horizontal arrows have ramification index $1$
and induce finite separable residue field extensions.
Thus, using the equivalence of
Lemma \href{https://stacks.math.columbia.edu/tag/09E7}{lemma extension dvrs formally smooth (Tag 09E7)},
we see that it suffices to show that the right vertical
arrow is formally smooth in the $\mathfrak m'_{ij}$-adic topology.
Since $u$ has a $e$th root
in $K_2$ we obtain the claim.

\medskip\noindent
Assume $u$ has an $e$th root in $K_1$.
Since $e | e_1$ and since $u$ has a $e$th root in $K_1$
we see that $\pi = \theta^e$ for some $\theta \in K_1$.
Let $K'_1 = K[\theta] \subset K_1$ be the subfield generated by $\theta$.
By Lemma \href{https://stacks.math.columbia.edu/tag/09EV}{lemma pull root uniformizer (Tag 09EV)} the integral closure $A'_1$
of $A$ in $K[\theta]$ is the discrete valuation ring $A'_1 = A[\theta]$
which has ramification index $e$ over $A$.
If we can prove the lemma for the extension $K'_1/K$,
then we conclude by Lemma \href{https://stacks.math.columbia.edu/tag/09EQ}{lemma formally smooth goes up (Tag 09EQ)}
applied to the diagram
$$
\xymatrix{
(B'_1)_{B'_1 \cap \mathfrak m_{ij}} \ar[r] & (B_1)_{\mathfrak m_{ij}} \\
A'_1 \ar[u] \ar[r] &
(A_1)_{\mathfrak m_i} \ar[u]
}
$$
for all $j = 1, \ldots, m_i$. This reduces us to the case discussed
in the next paragraph.

\medskip\noindent
Assume $K_1 = K[\pi^{1/e}]$ and set $\theta = \pi^{1/e}$. Let $\pi_B$ be a
uniformizer for $B$ and write $\pi = w \pi_B^e$ for some unit $w$ of $B$.
Then we see that $L_1 = L \otimes_K K_1$ is obtained by adjoining
$\theta / \pi_B$ which is an $e$th root of the unit $w$. Thus
$B \subset B_1$ is finite \'etale. Thus for any maximal ideal
$\mathfrak m \subset B_1$ consider the commutative diagram
$$
\xymatrix{
B \ar[r]_1 & (B_1)_{\mathfrak m} \\
A \ar[u]^e \ar[r]^e & A_1 \ar[u]_{e_\mathfrak m}
}
$$
Here the numbers along the arrows are the ramification indices.
By multiplicativity of ramification indices
(Lemma \href{https://stacks.math.columbia.edu/tag/0BRL}{lemma multiplicative e f (Tag 0BRL)})
we conclude $e_\mathfrak m = 1$. Looking at the residue field extensions
we find that $\kappa(\mathfrak m)$ is a finite separable extension
of $\kappa_B$ which is separable over $\kappa_A$. Therefore
$\kappa(\mathfrak m)$ is separable over $\kappa_A$
which is equal to the residue field of $A_1$ and we win by
Lemma \href{https://stacks.math.columbia.edu/tag/09E7}{lemma extension dvrs formally smooth (Tag 09E7)}.
\end{proof}
\end{document}
